\documentclass [10pt,letterpaper]{report}
%\documentclass[10pt,letterpaper]{article}
%\documentclass[10pt,letterpaper,twocolumn]{article}
%
%\documentclass[10pt,letterpaper,twocolumn]{report}
%\documentclass[10pt,letterpaper]{report}
%
%\usepackage[greek,english]{babel}
\usepackage[greek,english]{babel}
\usepackage[T1]{fontenc}
\usepackage{graphicx}
\usepackage{color}
\usepackage{wrapfig}


\textheight     =  8.4in
\textwidth      =  6.8in
%\parindent      =  0.0cm

\topmargin      =  -0.2in
\oddsidemargin  =  -0.2in
\evensidemargin =  -0.2in
%\hoffset        =  0.0cm
%\voffset        =  0.0cm

\newcommand{\nuSolve}{\greektext n\latintext Solve}
\newcommand{\thisVersion}{0.1.2}

\graphicspath{{graphics/}}

\begin{document}

\title{log2ant-\thisVersion{}: User Guide}

\author{
Sergei Bolotin,
Karen Baver,
John Gipson,
David Gordon,
Daniel MacMillan
}

\date{\today}


\maketitle

\tableofcontents

\def\leftSp      {48mm}
\def\riteSp     {113mm}



\chapter{Introduction}
\label{intro}

This document describes how to use the utility log2ant.

The utility performs extracting various sensors reading from a station log
file and store them in an ASCII file.

The utility log2ant is distributed in nusolve package that contains \nuSolve{}
software and utilities vgosDbMake, vgosDbCalc, vgosDbProcLogs and log2ant. Each of the
utilities has its own version number that can differ from distribution version,
e.g., the first release of log2ant-0.0.1 appeared in nusolve-0.7.2 package.

The guide covers \thisVersion{} version of the software. Since log2ant
is a simple utility we do not expect large modifications of the
user guide version from version.

\section{Requirements}
\label{intro_reqs}

See \emph{\nuSolve{} User Guide} for the requirements.


\section{Changes from previous versions}
\label{intro_changes}

This section covers changes in the software and the user guide from version to version.


\subsection{Changes in version 0.1.3}
Support of ANTCAL version 2 format was added.

\subsection{Changes in version 0.1.2}
Nothing essential was changed.

\subsection{Changes in version 0.1.1}
Nothing essential was changed.

\subsection{Changes in version 0.1.0}
Regular expressions to parse modern log files were updated. An empty string filler was
replaced from "n/a" to "DUMMY".
Values of dot2pps records are extracted and reported in an ANTCAL file.
Frequencies of PCAL sensors are calculated and reported in an ANTCAL file.

\subsection{Changes in version 0.0.5}
Nothing essential was changed.

\subsection{Changes in version 0.0.4}
Nothing essential was changed.

\subsection{Changes in version 0.0.3}

The command line arguments parser has switched to ARGP from GNU C Library.

Processing of log files was improved, the regular expressions were updated.


\subsection{Changes in version 0.0.2}
An ability to extract SEFD evaluation from log files is added. SEFD is calculated
by Field System as a result of observing a selected source (using procedure {\it onoff}).
Also, processing DBBC3 dump files is implemented (thanks to Eskil Varenius, Onsala Space Observatory).

A command line option to redirect utility's log output is added.
Section \emph{\ref{section_invokation} Invoking log2ant} was modified to reflect
these changes.

The utility can process a log where channel setup is changing. That could be useful
to process log files of single dish observations.

Support for LOG-ANTAB output format is suspended.

A few bugs were fixed.

\subsection{Changes in version 0.0.1}
Initial release.


%%%%%%%%%%%%%%%%%%%%%%%%%%%%%%%%%%%%%%%%%%%%%%%%%%%%%%%%%%%%%%%%%%%%%%%%%%%%%%
%
%
\chapter{Installation}

The source codes of log2ant is distributed along with \nuSolve{}
software. The latest stable version of the software one can find at
{\tt https://sourceforge.net/projects/nusolve} with a name like
{\tt nusolve-1.2.3.tar.gz}. Since the software is still in an active
development phase, we recommend you use the latest version.

The utility is compiled during compilation of \nuSolve{} software. Please
refer to \nuSolve{} User Guide how to configure, compile and install the
software.




%%%%%%%%%%%%%%%%%%%%%%%%%%%%%%%%%%%%%%%%%%%%%%%%%%%%%%%%%%%%%%%%%%%%%%%%%%%%%%
%
%
\chapter{Using log2ant}
\section{Invoking log2ant}
\label{section_invokation}
To invoke log2ant just type (specifying if necessary the full path
to the executable) program name and a file name of a field system log of a station:
\begin{verbatim}
> log2ant <FS log file>
\end{verbatim}
\noindent
where <<FS log file>> is a name of input file. The input file can be compressed with
gzip or bzip2, the utility will automatically decompress the file if gzip or bzip2
compressors are installed on you system and are in your {\tt \$PATH}. It is expected
that the type of compressed file is determined by the extension. If the file name
ends with <<.gz>>, log2ant assumes it was compressed with gzip, if it ends with
<<.bz2>>, bzip2 compressor will be used.

The utility also accepts command line arguments. To get the list of these
arguments, type

\begin{verbatim}
> log2ant --help
\end{verbatim}
\noindent
Here are command line arguments that are available at the time of writing:
\vspace{5mm}

\begin{tabular}{p{\leftSp}p{\riteSp}}
\multicolumn{2}{c}{General options:}\\
{\tt -s, -{}-station-name=STRING}  & Set a name of a station to STRING.\\
\end{tabular}
\vspace{5mm}

\begin{tabular}{p{\leftSp}p{\riteSp}}
\multicolumn{2}{c}{Input control:}\\
{\tt -3, -{}-DBBC3-dump=STRING}    & Use a file STRING as a DBBC3 dump file.\\
\end{tabular}
\vspace{5mm}

\begin{tabular}{p{\leftSp}p{\riteSp}}
\multicolumn{2}{c}{Output control:}\\
{\tt -c, -{}-compress=STRING}      & Use a compressor STRING (gzip or bzip2) to squeeze output ANTCAL file.\\
{\tt -o, -{}-output=STRING}        & Set a name of output ANTCAL file to STRING.\\
\end{tabular}
\vspace{5mm}

\begin{tabular}{p{\leftSp}p{\riteSp}}
\multicolumn{2}{c}{Time interval options:}\\
{\tt -b, -{}-t-begin=STRING}       & Set an epoch of the first observation to STRING, data before STRING will be ignored.\\
{\tt -e, -{}-t-end=STRING}         & Set an epoch of the last observation to STRING, data after STRING will be ignored.\\
\end{tabular}
\vspace{5mm}

\begin{tabular}{p{\leftSp}p{\riteSp}}
\multicolumn{2}{c}{Data filter:}\\
{\tt -a, -{}-all}                  & Report all collected data (i.e., both data\_valid on and off intervals).\\
{\tt -t, -{}-data-type=STRING}     & Extract only the specified type of data. STRING can be: cbl (cable calibration), dat
(data=on/off), fmt (fmt2gps), met (meteorological parameters), phc (phase calibration), sefd (SEFD evaluation),
tpc (DBBC3 TPC), tpi (TPI), tsys (TSYS). There can be more than one "-t" option, e.g.: -t dat -t tsys.\\
{\tt -u, -{}-strip-unused-sensors} & Do not output sensors that are not in a log file for some particular time.\\
\end{tabular}
\vspace{5mm}

\begin{tabular}{p{\leftSp}p{\riteSp}}
\multicolumn{2}{c}{Debug output:}\\
{\tt -l, -{}-log-file=STRING}      & Store log2ant's output in a file STRING, the default is "log2ant.log".\\
{\tt -v, -{}-verbose=NUM}          & Set a level of log output to NUM: 0 (errors only), 1 (+warnings), 2 (+info) and 3 (+debug output). Default is 1.\\
\end{tabular}
\vspace{5mm}

\begin{tabular}{p{\leftSp}p{\riteSp}}
\multicolumn{2}{c}{Operation modes:}\\
{\tt -?, -{}-help   }              & Give this help list.\\
{\tt     -{}-usage  }              & Give a short usage message.\\
{\tt -V, -{}-version}              & Print program version.\\
\end{tabular}
\vspace{5mm}

\noindent
The options expect the following meaning.

The option <<-3 DBBC3.dump>> specifies a file name of the corresponding DBBC3 dump file.
The software parses such a file when it is available along with the FS log file. For
example,
\begin{verbatim}
> log2ant -3 vo0202oe.dump.gz vo0202oe.log.gz
\end{verbatim}
\noindent
will process the Field System log file {\tt vo0202oe.log.gz} and extract data from the DBBC3
dump file {\tt vo0202oe.dump.gz}.

The option <<-a>> tells log2ant to collect all
readings from a log file. Without this option it will collect data only between
data\_valid=on and data\_valid=off records in a log file.

The option <<-b tBegin>> is necessary when a log file contains records from several
sessions. When log2ant parses a log file, it has this information from a database.
The utility log2ant has no such information, so in some cases it is useful to provide
epoch of start of the session. It expects an epoch of first observations, tBegin, in the
same format as it was used in a log file. E.g.:
\begin{verbatim}
> log2ant -b 2020.021.17:40:16.09 /500/sessions/2020/vo0021/vo0021gs.log
\end{verbatim}

\noindent
also, a user can use more human notation providing year, month and day instead of
Field System format:

\begin{verbatim}
> log2ant -b 2020.01.21.17:40:16.09 /500/sessions/2020/vo0021/vo0021gs.log
\end{verbatim}

The option <<-c>> turns on using a compressor for the output file. Currently, two
compressors are supported: gzip (<<-c gz>>) and bzip2 (<<-c bz2>>). In addition, if
a user specifies output file name and the name ends with <<.gz>> or <<.bz2>>,
the corresponding compressor will be used automatically and the option <<-c>> is
not necessary.

The option <<-e tEnd>> is the similar to <<-b>> option and sets the end of the session.
Usually, after the last observation of a session there are measurements of cable length
with and without additional cable to determine a cable sign. So, the utility reads records
form a log file up to tEnd + 5 minutes.

%The option <<-f>> controls format of output file. Currently, two types are implemented:
%ANTCAL (default) and LOG-ANTAB.

The option <<-o>> specifies a file name of the output. By default, log2ant will try
to write the file in the same directory where the log file is, modifying name:
".log" is replaced with ".anc" and underscore char, "\_", is added before last two
chars of the name. E.g., if input file name is "vo0021gs.log", then the default
output file name will be "vo0021\_gs.anc".

Sometimes a station name in a log file is not correct or not present at all. In such cases
a user can specify station name using option <<-s stationName>>.

The option <<-t dataType>> controls what type of data should be written in the output file.
The following types are known to the utility:

\begin{tabular}{ll}
{\tt  cbl}  & cable calibrations (converted to seconds)\\
{\tt  dat}  & data\_valid=on records\\
{\tt  fmt}  & fmt2gps readings\\
{\tt  met}  & meteorological parameters\\
{\tt  phc}  & phase calibrations\\
{\tt sefd}  & evaluations of SEFD\\
{\tt  tpc}  & DBBC3 tpc readings\\
{\tt  tpi}  & tpi readings\\
{\tt tsys}  & calculated by FS tsys values\\
\end{tabular}

\noindent
Each of type data can be combined with another one. For example:

\begin{verbatim}
> log2ant -t cbl -t met /500/sessions/2020/vo0021/vo0021gs.log
\end{verbatim}

\noindent
will extract only cable calibrations and meteorological parameters. If this option is not
provided, all found data will be reported.

The option <<-v>> tunes the level of log output. The possible levels are: error (0),
warning (1), informational (2) and debug (3) messages. Currently by default the log level
is set to 1, that means that only error and warning messages will be displayed.


\section{ANTCAL format}
\label{section_antcal_format}

The format ANTCAL is currently in a developing stage. In future it could be extended.
Current version of the format is available in nusolve distribution in 
{\tt docs/antcal\_specs\_???.txt}.




%%%%%%%%%%%%%%%%%%%%%%%%%%%%%%%%%%%%%%%%%%%%%%%%%%%%%%%%%%%%%%%%%%%%%%%%%%%%%%
%
%
\chapter{Concluding remark}
\label{chapt_concluding_remark}
Currently, this document is in the developmental stage, its content could change
time from time. Check for new versions at the ftp site:

\begin{verbatim}
      https://sourceforge.net/projects/nusolve
\end{verbatim}

\noindent
If you have questions or suggestions that will improve the software or
the User Guide, please e-mail us at:

\begin{verbatim}
      <mailto:sergei.bolotin@nasa.gov>
\end{verbatim}

%\noindent
%Happy !


%%%%%%%%%%%%%%%%%%%%%%%%%%%%%%%%%%%%%%%%%%%%%%%%%%%%%%%%%%%%%%%%%%%%%%%%%%%%%%%
%%
%%
%\begin{thebibliography}{99}
%
%\bibitem{Bolotin2010}
%S. Bolotin, J. Gipson, D. MacMillan: ``Development of a New VLBI Data
%Analysis Software". In: International VLBI Service for Geodesy and
%Astrometry 2010 General Meeting Proceedings, edited by D. Behrend and K.
%Baver, NASA/CP-2010-215864, 197-201, 2010.
%
%
%\bibitem{Bolotin2011}
%S. Bolotin, J. Gipson, D. Gordon, D. MacMillan: ``Current Status of
%Development of New VLBI Data Analysis Software". In: Proceedings of the
%20th Meeting of the European VLBI Group for Geodesy and Astrometry,
%edited by W. Alef, S. Bernhart, A. Nothnagel, Schriftenreihe des
%Instituts f\"ur Geod\"asie und Geoinformation der Universit\"at Bonn,
%Nr. 22, ISSN 1864-1113, 86-88, 2011.
%
%
%\bibitem{Bolotin2012}
%S. Bolotin, K. Baver, J. Gipson, D. Gordon, D. MacMillan: ``The First
%Release of $\nu$Solve". In: International VLBI Service for Geodesy and
%Astrometry 2012 General Meeting Proceedings `Launching the
%Next-Generation IVS Network', edited by D. Behrend and K. Baver,
%NASA/CP-2012-217504, 222-226, 2012.
%
%
%\end{thebibliography}


\end{document}


